\documentclass[11pt,a4paper]{article}

\usepackage[T1]{fontenc}
\usepackage[utf8]{inputenc}
\usepackage[french]{babel}
\usepackage[margin=2.5cm]{geometry}
\usepackage{booktabs}
\usepackage{siunitx}
\usepackage[table]{xcolor}
\usepackage{hyperref}
\usepackage{longtable}
\usepackage{pgfplots}
\pgfplotsset{
  compat=1.18,
  every axis legend/.append style={
    at={(0.5,0.98)}, anchor=north,
    legend columns=4,
    font=\scriptsize,
    fill=white, fill opacity=0.85, draw opacity=1, text opacity=1,
    draw=black!20,
    /tikz/every even column/.append style={column sep=6pt},
  },
}
\usepackage{caption}
\usepackage{enumitem}
\usepackage{amsmath}
\usepackage{amssymb}

\sisetup{detect-weight=true, detect-family=true}

\hypersetup{
  colorlinks=true,
  linkcolor=blue!50!black,
  urlcolor=blue!50!black,
  citecolor=blue!50!black,
  pdftitle={Rapport d'activite - Scan H(P) r2SCAN Al4C3 (partiel)},
  pdfauthor={Pr. John Akapulco}
}

\title{\vspace{-1.5cm}Rapport d'activité\\[2pt]
\large Diagramme de stabilité H(P) r2SCAN des polymorphes d'Al$_4$C$_3$\\[2pt]
\normalsize (résultats partiels, référence R-3m)}
\author{Pr. John Akapulco \\ \small \href{mailto:computchem86@gmail.com}{computchem86@gmail.com}}
\date{23 août 2026}

\begin{document}
\maketitle
\vspace{-0.5cm}

\begin{abstract}
\noindent
Ce rapport présente l'état d'avancement du scan H(P) r2SCAN (huit polymorphes, grille 0 à
100~GPa par pas de 10~GPa, méthodologie décrite dans
\texttt{rapport\_activite\_Al4C3\_PBE\_vs\_r2SCAN.tex}) au 23~août~2026 (soir) :
\textbf{66 points convergés sur 88}. Cinq polymorphes sont complets (Fd-3m, Pnma-VII,
Pnma-III, Pnma, R-3m) ; trois sont encore en cours -- \textbf{Pbam} (6/11), \textbf{I-43d}
(4/11, chaîne relancée à P40 après un plantage numérique, job SLURM 58469) et
\textbf{anti-CaFe$_2$O$_4$} (1/11, job SLURM 58468). Les diagrammes $V(P)$, $E(P)$ et $H(P)$
ci-dessous n'incluent donc que les points déjà convergés ; les courbes des trois polymorphes
incomplets s'arrêtent à la dernière pression atteinte. Le diagramme d'enthalpie est présenté en
valeur \emph{relative}, $\Delta H(P) = H_{\text{phase}}(P) - H_{\text{R-3m}}(P)$, avec
\textbf{R-3m} comme référence (même convention que
\texttt{note\_Al4C3\_enthalpie\_relative\_v2.tex} pour le jeu PBE), R-3m étant complet sur toute
la gamme et servant donc de référence commune fiable même pour les polymorphes partiellement
convergés.
\end{abstract}

\section{État de convergence}

\begin{center}
\begin{tabular}{lc}
\toprule
Polymorphe & Points convergés (0--100~GPa) \\
\midrule
Fd-3m              & 11/11 \\
Pnma-VII           & 11/11 \\
Pnma-III           & 11/11 \\
Pnma               & 11/11 \\
R-3m               & 11/11 \\
Pbam               & 6/11 (0--50~GPa) \\
I-43d              & 4/11 (0--30~GPa) \\
anti-CaFe$_2$O$_4$ & 1/11 (0~GPa) \\
\midrule
\textbf{Total} & \textbf{66/88} \\
\bottomrule
\end{tabular}
\label{tab:r2scan_conv}
\end{center}

Convergence = critère de forces \texttt{EDIFFG=-0.01}~eV/\AA\ atteint (\emph{"reached required
accuracy"} dans l'\texttt{OUTCAR}). Pbam et anti-CaFe$_2$O$_4$ tournent normalement (jobs SLURM
58467/58468). La chaîne I-43d avait planté à P40 (\texttt{Error EDDDAV: Call to ZHEGV failed},
instabilité numérique de la diagonalisation itérative, sans lien avec la qualité de la seed
PBE) ; elle a été nettoyée et relancée de P40 à P100 (job 58469).

\section{Diagramme de volume V(P)}

\begin{figure}[h]
\centering
\begin{tikzpicture}
\begin{axis}[
  width=\textwidth, height=8cm,
  xlabel={Pression (GPa)},
  ylabel={$V$/formule (\AA$^3$)},
  enlarge y limits={upper=0.22, lower=0.02},
  grid=major,
]
\addplot table[x=P_GPa,y=V_A3_per_fu] {../../hp_curves_r2SCAN/hp_V_Pnma-VII.dat};
\addplot table[x=P_GPa,y=V_A3_per_fu] {../../hp_curves_r2SCAN/hp_V_Pbam.dat};
\addplot table[x=P_GPa,y=V_A3_per_fu] {../../hp_curves_r2SCAN/hp_V_Pnma.dat};
\addplot table[x=P_GPa,y=V_A3_per_fu] {../../hp_curves_r2SCAN/hp_V_R-3m.dat};
\addplot table[x=P_GPa,y=V_A3_per_fu] {../../hp_curves_r2SCAN/hp_V_Pnma-III.dat};
\addplot table[x=P_GPa,y=V_A3_per_fu] {../../hp_curves_r2SCAN/hp_V_anti-CaFe2O4.dat};
\addplot table[x=P_GPa,y=V_A3_per_fu] {../../hp_curves_r2SCAN/hp_V_I-43d.dat};
\addplot table[x=P_GPa,y=V_A3_per_fu] {../../hp_curves_r2SCAN/hp_V_Fd-3m.dat};
\legend{Pnma-VII, Pbam, Pnma, R-3m, Pnma-III, anti-CaFe2O4, I-43d, Fd-3m}
\end{axis}
\end{tikzpicture}
\caption{Volume de la maille par formule unitaire (Al$_4$C$_3$), r2SCAN, en fonction de la
pression. Les courbes Pbam, I-43d et anti-CaFe$_2$O$_4$ s'arrêtent à leur dernier point
convergé (respectivement 50, 30 et 0~GPa) -- scan encore en cours pour ces trois polymorphes.}
\label{fig:vp_r2scan}
\end{figure}

\section{Diagramme d'énergie totale E(P)}

\begin{figure}[h]
\centering
\begin{tikzpicture}
\begin{axis}[
  width=\textwidth, height=8cm,
  xlabel={Pression (GPa)},
  ylabel={$E$/formule (eV)},
  enlarge y limits={upper=0.22, lower=0.02},
  grid=major,
]
\addplot table[x=P_GPa,y=E_eV_per_fu] {../../hp_curves_r2SCAN/hp_E_Pnma-VII.dat};
\addplot table[x=P_GPa,y=E_eV_per_fu] {../../hp_curves_r2SCAN/hp_E_Pbam.dat};
\addplot table[x=P_GPa,y=E_eV_per_fu] {../../hp_curves_r2SCAN/hp_E_Pnma.dat};
\addplot table[x=P_GPa,y=E_eV_per_fu] {../../hp_curves_r2SCAN/hp_E_R-3m.dat};
\addplot table[x=P_GPa,y=E_eV_per_fu] {../../hp_curves_r2SCAN/hp_E_Pnma-III.dat};
\addplot table[x=P_GPa,y=E_eV_per_fu] {../../hp_curves_r2SCAN/hp_E_anti-CaFe2O4.dat};
\addplot table[x=P_GPa,y=E_eV_per_fu] {../../hp_curves_r2SCAN/hp_E_I-43d.dat};
\addplot table[x=P_GPa,y=E_eV_per_fu] {../../hp_curves_r2SCAN/hp_E_Fd-3m.dat};
\legend{Pnma-VII, Pbam, Pnma, R-3m, Pnma-III, anti-CaFe2O4, I-43d, Fd-3m}
\end{axis}
\end{tikzpicture}
\caption{Énergie totale (sans le terme $PV$) par formule unitaire (Al$_4$C$_3$), r2SCAN, en
fonction de la pression, mêmes points convergés que la figure~\ref{fig:vp_r2scan}.}
\label{fig:ep_r2scan}
\end{figure}

\section{Diagramme d'enthalpie relative $\Delta H(P)$, référence R-3m}

\begin{figure}[h]
\centering
\begin{tikzpicture}
\begin{axis}[
  width=\textwidth, height=8.5cm,
  xlabel={Pression (GPa)},
  ylabel={$\Delta H$/atome (meV), réf. R-3m},
  enlarge y limits={upper=0.22, lower=0.02},
  grid=major,
  extra y ticks={0},
  extra y tick style={grid style={solid, black!40}},
]
\addplot table[x=P_GPa,y=dH_meV_per_atom] {../../hp_curves_r2SCAN/hp_rel_R-3m.dat};
\addplot table[x=P_GPa,y=dH_meV_per_atom] {../../hp_curves_r2SCAN/hp_rel_Pnma-VII.dat};
\addplot table[x=P_GPa,y=dH_meV_per_atom] {../../hp_curves_r2SCAN/hp_rel_Fd-3m.dat};
\addplot table[x=P_GPa,y=dH_meV_per_atom] {../../hp_curves_r2SCAN/hp_rel_Pbam.dat};
\addplot table[x=P_GPa,y=dH_meV_per_atom] {../../hp_curves_r2SCAN/hp_rel_Pnma-III.dat};
\addplot table[x=P_GPa,y=dH_meV_per_atom] {../../hp_curves_r2SCAN/hp_rel_Pnma.dat};
\addplot table[x=P_GPa,y=dH_meV_per_atom] {../../hp_curves_r2SCAN/hp_rel_anti-CaFe2O4.dat};
\addplot table[x=P_GPa,y=dH_meV_per_atom] {../../hp_curves_r2SCAN/hp_rel_I-43d.dat};
\legend{R-3m (réf.), Pnma-VII, Fd-3m, Pbam, Pnma-III, Pnma, anti-CaFe2O4, I-43d}
\end{axis}
\end{tikzpicture}
\caption{Enthalpie relative $\Delta H(P) = H_{\text{phase}}(P) - H_{\text{R-3m}}(P)$, r2SCAN, en
meV/atome. R-3m est la droite $\Delta H = 0$. anti-CaFe$_2$O$_4$ n'a qu'un point (0~GPa) ; Pbam
et I-43d s'arrêtent respectivement à 50 et 30~GPa. Comme pour le jeu PBE
(\texttt{note\_Al4C3\_enthalpie\_relative\_v2.tex}), Pnma-VII reste sous zéro et Fd-3m la
rejoint sous zéro à basse pression ; les autres phases démarrent au-dessus de zéro à 0~GPa.}
\label{fig:hp_rel_r2scan}
\end{figure}

\section{Valeurs numériques (E, V, H) -- points convergés à ce jour}

Le tableau~\ref{tab:hp_r2scan_full} rassemble, pour les 66 points convergés, l'énergie totale
$E$, le volume $V$ et l'enthalpie $H$ par formule unitaire (r2SCAN), arrondis au millième.

% Auto-generated from hp_curves_r2SCAN/results_hp.csv -- do not edit by hand.
% Regenerate with: python3 hp_curves_r2SCAN/extract_hp.py, then re-run the
% table-export snippet noted in the report's methodology section.
\begin{longtable}{lrS[table-format=-2.3]S[table-format=2.3]S[table-format=-2.3]}
\caption{Énergie totale $E$, volume $V$ et enthalpie $H = E + PV$ par formule unitaire
(Al$_4$C$_3$, 7 atomes), r2SCAN, points convergés à ce jour (66/88 -- voir
tableau~\ref{tab:r2scan_conv} pour l'état par polymorphe). Valeurs arrondies au millième
d'eV (resp.\ d'\AA$^3$).}
\label{tab:hp_r2scan_full} \\
\toprule
Polymorphe & {$P$ (GPa)} & {$E$ (eV)} & {$V$ (\AA$^3$)} & {$H$ (eV)} \\
\midrule
\endfirsthead
\multicolumn{5}{l}{\small \textit{(suite)}} \\
\toprule
Polymorphe & {$P$ (GPa)} & {$E$ (eV)} & {$V$ (\AA$^3$)} & {$H$ (eV)} \\
\midrule
\endhead
\bottomrule
\endfoot
Pnma-VII & 0 & -58.458 & 79.733 & -58.458 \\
Pnma-VII & 10 & -58.332 & 75.638 & -53.611 \\
Pnma-VII & 20 & -58.025 & 72.373 & -48.991 \\
Pnma-VII & 30 & -57.601 & 69.657 & -44.557 \\
Pnma-VII & 40 & -57.093 & 67.340 & -40.281 \\
Pnma-VII & 50 & -56.524 & 65.323 & -36.139 \\
Pnma-VII & 60 & -55.909 & 63.532 & -32.117 \\
Pnma-VII & 70 & -55.263 & 61.943 & -28.200 \\
Pnma-VII & 80 & -54.580 & 60.483 & -24.379 \\
Pnma-VII & 90 & -53.877 & 59.160 & -20.644 \\
Pnma-VII & 100 & -53.156 & 57.945 & -16.989 \\
\midrule
Pbam & 0 & -57.151 & 71.880 & -57.151 \\
Pbam & 10 & -57.044 & 68.430 & -52.773 \\
Pbam & 20 & -56.785 & 65.670 & -48.587 \\
Pbam & 30 & -56.424 & 63.365 & -44.559 \\
Pbam & 40 & -55.990 & 61.390 & -40.664 \\
Pbam & 50 & -55.504 & 59.665 & -36.884 \\
\midrule
Pnma & 0 & -56.471 & 69.912 & -56.471 \\
Pnma & 10 & -56.364 & 66.415 & -52.218 \\
Pnma & 20 & -56.104 & 63.640 & -48.159 \\
Pnma & 30 & -55.748 & 61.367 & -44.257 \\
Pnma & 40 & -55.323 & 59.430 & -40.486 \\
Pnma & 50 & -54.847 & 57.740 & -36.828 \\
Pnma & 60 & -54.335 & 56.252 & -33.269 \\
Pnma & 70 & -53.791 & 54.915 & -29.798 \\
Pnma & 80 & -53.223 & 53.708 & -26.407 \\
Pnma & 90 & -52.636 & 52.602 & -23.087 \\
Pnma & 100 & -52.037 & 51.595 & -19.835 \\
\midrule
R-3m & 0 & -58.432 & 79.937 & -58.432 \\
R-3m & 10 & -58.306 & 75.860 & -53.571 \\
R-3m & 20 & -58.000 & 72.597 & -48.938 \\
R-3m & 30 & -57.577 & 69.883 & -44.491 \\
R-3m & 40 & -57.069 & 67.567 & -40.200 \\
R-3m & 50 & -56.500 & 65.543 & -36.045 \\
R-3m & 60 & -55.885 & 63.757 & -32.009 \\
R-3m & 70 & -55.234 & 62.153 & -28.078 \\
R-3m & 80 & -54.554 & 60.703 & -24.243 \\
R-3m & 90 & -53.852 & 59.383 & -20.494 \\
R-3m & 100 & -53.131 & 58.170 & -16.824 \\
\midrule
Pnma-III & 0 & -57.315 & 72.853 & -57.315 \\
Pnma-III & 10 & -57.201 & 69.185 & -52.883 \\
Pnma-III & 20 & -56.929 & 66.282 & -48.655 \\
Pnma-III & 30 & -56.552 & 63.870 & -44.593 \\
Pnma-III & 40 & -56.101 & 61.810 & -40.669 \\
Pnma-III & 50 & -55.595 & 60.015 & -36.866 \\
Pnma-III & 60 & -55.049 & 58.428 & -33.169 \\
Pnma-III & 70 & -54.470 & 57.005 & -29.565 \\
Pnma-III & 80 & -53.866 & 55.715 & -26.046 \\
Pnma-III & 90 & -53.242 & 54.540 & -22.604 \\
Pnma-III & 100 & -52.600 & 53.462 & -19.232 \\
\midrule
anti-CaFe$_2$O$_4$ & 0 & -56.363 & 70.440 & -56.363 \\
\midrule
I-43d & 0 & -54.982 & 68.155 & -54.982 \\
I-43d & 10 & -54.863 & 64.155 & -50.859 \\
I-43d & 20 & -54.605 & 61.388 & -46.942 \\
I-43d & 30 & -54.259 & 59.172 & -43.179 \\
\midrule
Fd-3m & 0 & -58.340 & 77.912 & -58.340 \\
Fd-3m & 10 & -58.222 & 74.000 & -53.603 \\
Fd-3m & 20 & -57.927 & 70.865 & -49.081 \\
Fd-3m & 30 & -57.517 & 68.251 & -44.737 \\
Fd-3m & 40 & -57.026 & 66.014 & -40.545 \\
Fd-3m & 50 & -56.476 & 64.062 & -36.484 \\
Fd-3m & 60 & -55.880 & 62.333 & -32.536 \\
Fd-3m & 70 & -55.249 & 60.784 & -28.693 \\
Fd-3m & 80 & -54.592 & 59.381 & -24.942 \\
Fd-3m & 90 & -53.912 & 58.101 & -21.274 \\
Fd-3m & 100 & -53.214 & 56.925 & -17.684 \\

\end{longtable}


\appendix
\section{Références du dépôt}

Données brutes dans \texttt{hp\_curves\_r2SCAN/Al4C3\_<polymorphe>/P<pression>/}, extraction via
\texttt{hp\_curves\_r2SCAN/extract\_hp.py} (\texttt{results\_hp.csv},
\texttt{hp\_<polymorphe>.dat} (H), \texttt{hp\_E\_<polymorphe>.dat} (E),
\texttt{hp\_V\_<polymorphe>.dat} (V)). Fichiers relatifs (référence R-3m, meV/atome) :
\texttt{hp\_curves\_r2SCAN/hp\_rel\_<polymorphe>.dat}. Tableau complet (section~5) dans
\texttt{table\_hp\_r2scan\_full.tex} / \texttt{table\_hp\_r2scan\_full\_body.tex}. Chaîne I-43d
relancée : \texttt{hp\_curves\_r2SCAN/Al4C3\_I-43d/run\_chain\_resume\_P40.sh} (job SLURM
58469). Voir \texttt{rapport\_activite\_Al4C3\_HP\_scan.tex} pour le scan PBE complet et
\texttt{rapport\_activite\_Al4C3\_PBE\_vs\_r2SCAN.tex} pour la méthodologie DFT r2SCAN
(seed PBE, \texttt{METAGGA=R2scan}, etc.).

\end{document}
